\documentclass[%
 reprint,
%superscriptaddress,
%groupedaddress,
%unsortedaddress,
%runinaddress,
%frontmatterverbose, 
%preprint,
%preprintnumbers,
%nofootinbib,
%nobibnotes,
%bibnotes,
 amsmath,amssymb,
 %aip,
 aps,
%pra,
%prb,
%rmp,
%prstab,
%prstper,
%floatfix,
]{revtex4-2}

\usepackage{graphicx}% Include figure files
%\graphicspath{{./Figs}}
%\usepackage{dcolumn}% Align table columns on decimal point
%\usepackage{bm}% bold math
%\newcommand{\bvec}[1]{\bm{\mathrm{#1}}}
%\usepackage{xcolor,soul}
%\usepackage[ISO]{diffcoeff}
%\usepackage{chemformula}

\usepackage{hyperref}% add hypertext capabilities
% From CUED PhD Latex template class definition
%\hypersetup{
 %   final=true,
  %  plainpages=false,
   % pdfstartview=FitV,
    %pdftoolbar=true,
    %pdfmenubar=true,
    %bookmarksopen=true,
    %bookmarksnumbered=true,
    %breaklinks=true,
    %linktocpage,
    %colorlinks=true,
    %linkcolor=blue,
    %urlcolor=blue,
    %citecolor=blue,
    %anchorcolor=black
  %}
%\usepackage[nameinlink]{cleveref}


%\usepackage{hyperref}% add hypertext capabilities
%\usepackage[mathlines]{lineno}% Enable numbering of text and display math
%\linenumbers\relax % Commence numbering lines

%\usepackage[showframe,%Uncomment any one of the following lines to test 
%%scale=0.7, marginratio={1:1, 2:3}, ignoreall,% default settings
%%text={7in,10in},centering,
%%margin=1.5in,
%%total={6.5in,8.75in}, top=1.2in, left=0.9in, includefoot,
%%height=10in,a5paper,hmargin={3cm,0.8in},
%]{geometry}
\usepackage{braket}
\usepackage{cleveref}
\Crefrangeformat{figure}{Fig.~#3#1#4 --~#5#2#6} 
\Crefrangeformat{table}{Tables~#3#1#4 --~#5#2#6} 
\Crefname{figure}{Fig.}{Fig.}
\Crefname{equation}{Eq.}{Eq.}
\creflabelformat{equation}{#2#1#3}
\usepackage{tikz}
\usetikzlibrary{matrix}
\usetikzlibrary{quantikz}
\usepackage{subfig}
\usepackage{todonotes}
\usepackage{ragged2e}
\usepackage{multirow}
\begin{document}

\title{Fighting noise with noise: a stochastic projective quantum eigensolver}
\author{Maria-Andreea Filip} \email{maf63@cam.ac.uk} 

%Change this if needed
\affiliation{ Yusuf Hamied Department of Chemistry, University of Cambridge,
Cambridge, UK }% 
\affiliation{ Peterhouse,
Cambridge, UK }% 
\date{\today}% It is always \today, today,
             %  but any date may be explicitly specified

\begin{abstract}
%Quantum Monte Carlo (QMC) algorithms have proven extremely effective at lowering the computational overhead of electronic structure calculations in a classical setting. In the current noisy intermediate scale quantum (NISQ) era of quantum computation, there are several limitations on the available hardware resources, such as low qubit counts, short decoherence times and high gate noise, which preclude the application of many current hybrid quantum-classical algorithms to non-trivial quantum chemistry problems. Here, we propose combining some of the fundamental elements of conventional QMC algorithms --- stochastic sampling of both the wavefunction and the Hamiltonian of interest --- with an imaginary-time propagation based projective quantum eigensolver. At the cost of increased noise, which can be easily averaged over in a classical Monte Carlo estimation, we obtain a method with quantum computational requirements that are both generally low and highly tunable.

In the current noisy intermediate scale quantum (NISQ) era of quantum computation, available hardware is severely limited by both qubit count and noise levels, precluding the application of many current hybrid quantum-classical algorithms to non-trivial quantum chemistry problems. In this paper we propose applying some of the fundamental ideas of conventional Quantum Monte Carlo algorithms --- stochastic sampling of both the wavefunction and the Hamiltonian --- to quantum algorithms in order to significantly decrease quantum resource costs. In the context of an imaginary-time propagation based projective quantum eigensolver, we present a novel approach to estimating physical observables which leads to an up to 50-fold reduction in the required sampling of the quantum state to converge the ground state energy of a system relative to current state-of-the-art eigensolvers. The method can be equally applied to excited-state calculations and, combined with stochastic truncations of the wavefunction Ansatz and system Hamiltonian, provides promising near-term approach to general chemistry Hamiltonian simulation on quantum devices.
\end{abstract}
\maketitle

\section{Introduction}
Electronic structure elucidation is a hard problem for conventional computation, due to the exponential scaling of the Hilbert space with the physical size of the system. Quantum computers have been heralded as a solution to this problem,\cite{Feynman1982, Aspuru-Guzik2005} mainly due to their ability to use the quantum entanglement of their constituent qubits to encode even highly complex electronic wavefunctions in a linearly scaling qubit register. This promise has led to the development of a variety of quantum algorithms with applications in quantum chemistry, such as Hamiltonian simulation\cite{Berry2007, Berry2015, Low2017, Childs2018} and Quantum Phase Estimation.\cite{Abrams1997, Abrams1999}

However, while these techniques promise clean, general solutions to many problems in quantum chemistry, they require many-qubit, fully fault-tolerant quantum computers to run. Currently we are in an era of Noisy Intermediate Scale Quantum (NISQ) hardware,\cite{Preskill2018} which is characterised by moderate numbers (10 -- 1000) of qubits prone to high error rates, due to short coherence times and noisy gate implementations.\cite{Linke2017} While significant work is being carried out in the fields of quantum error correction\cite{Georgescu2020} and mitigation, \cite{Cai2022} the fully fault-tolerant quantum computer remains a thing of the future.

In order to make the most of the currently available quantum resources, hybrid quantum-classical algorithms\cite{Farhi2014, Peruzzo2014, Colless2018, Higgott2019, Nakanishi2019, Greene-Diniz2021, Smart2021} have been developed, which shift most of the computational overhead onto a classical processor, while requiring specific, highly efficient operations to be carried out by the quantum processor. The most well known of these is the variational quantum eigensolver (VQE),\cite{Peruzzo2014} in which some parametrised wavefunction is encoded on a quantum device and used to compute a cost function --- usually the expectation value of a Hamiltonian operator. This is used in the classical optimisation of the parameters. In quantum chemistry, this approach has been used to calculate ground-state energies of molecules and model systems.\cite{Peruzzo2014, Shen2017, Hempel2018} 

While methods like this are able to make use of only moderate quantum resources, they still require a large number of repetitions of the necessary quantum measurements to overcome the underlying noise and obtain sufficiently accurate results to be of interest. 

In this paper, we take inspiration from classical Quantum Monte Carlo (QMC) algorithms,\cite{Booth2009, Thom2010, Filip2020} which use noisy representations of the wavefunction and stochastic propagation techniques to obtain very high-accuracy results for challenging systems. We use these ideas to decrease the quantum computational overhead for electronic structure problems. In particular, we consider a projective quantum eigensolver (PQE)\cite{Stair2021} approach and propose an algorithm to estimate the energy \textit{during} the propagation, which requires significantly fewer samples for the same accuracy as its conventional counterpart. We then introduce the idea of stochastic wavefunction rounding to reduce the depth of our Ansatz circuits and estimate the energy at each iteration using a sampled subset of the Hamiltonian, lowering the number of measurements required further. These techniques introduce further noise into the algorithm, but we find this can be efficiently averaged out and, for a set of second row hydrides, show that the accuracy of the original algorithm can be recovered at significantly lower cost. Finally, we investigate the applicability of this approach for excited states, using the folded-spectrum method, \cite{Peruzzo2014,Cao2019,CadiTazi2023} and find that, while convergence is somewhat more challenging, it can reduce the cost of an excited state calculation to that of the ground-state method.

In Sec. II we present the underlying theory of quantum eigensolvers, classical QMC methods and the folded-spectrum method which allows these algorithms to be expanded to excited state calculations. In Sec. III we describe the quantum algorithm employed in this paper, together with circuit implementations and possible stochastic variations. Finally, in Sec. IV we apply these methods to a range of molecular systems and discuss their performance. We draw our final conclusions in Sec. V.

\section{Theory}

\subsection{Variational Quantum Eigensolver}

One of the principal NISQ algorithms for quantum chemistry is the variational quantum eigensolver (VQE).\cite{Peruzzo2014} According to the variational principle, given a parametrised wavefunction $\Psi(\boldsymbol{\theta})$ that satisfies the boundary conditions of the problem, the expectation value of the energy is greater or equal to the ground state energy of the system,
\begin{equation}
\braket{E(\boldsymbol{\theta})} = \frac{\braket{\Psi(\boldsymbol\theta)|\hat H | \Psi(\boldsymbol{\theta})}}{\braket{\Psi(\boldsymbol{\theta})| \Psi(\boldsymbol{\theta})}} \geq E_0.
\end{equation}
Therefore, by minimising $\braket{E(\boldsymbol{\theta})}$ with respect to $\boldsymbol{\theta}$, one can obtain the best possible approximation of a given form to the ground state wavefunction.
In VQE, the expectation value of the energy and any gradients used are computed using a quantum processor before being passed on to a classical optimisation algorithm.

There are various considerations to be taken into account when choosing parametrised wavefunction Ans\"atze for VQE. Heuristic hardware efficient Ans\"atze (HEA)\cite{Kandala2017} attempt to introduce many degrees of freedom while maintaining low circuit depths. However, optimisation on the resulting landscapes is often difficult, as they are plagued by barren plateaus\cite{McClean2018} --- regions of parameter space in which the gradient vanishes exponentially with system size, which can prevent the optimisation algorithm from converging to a true minimum. This problem has been shown to be worse for more expressive Ans\"atze \cite{Holmes2022} and higher degrees of entanglement.\cite{Marrero2021} It has also been shown however that the intrinsic noise resulting from finite sampling of the quantum state can be helpful in improving optimisation outcomes.\cite{Liu2023} Alternatively, one can use physically motivated Ans\"atze such as unitary coupled cluster (UCC),\cite{Kutzelnigg1982,Kutzelnigg1983,Kutzelnigg1984,Bartlett1989} in which the wavefunction is expressed as
\begin{equation}
\ket{\Psi} = e^{\hat \tau}\ket{\phi_0} = e^{\hat T - \hat T^\dagger}\ket{\phi_0},
\end{equation}
where
\begin{align}
\hat T = \sum_\mathbf{i} t_\mathbf{i} \hat a_\mathbf{i},\\
\hat T^\dagger = \sum_\mathbf{i} t_\mathbf{i} \hat a_\mathbf{i}^\dagger
\end{align}
and $\hat a_\mathbf{i}$ and $\hat a_\mathbf{i}^\dagger$ are excitation and de-excitation operators respectively. This type of Ansatz is generally found to be easier to optimise than HEAs but has its own challenges. In order to express the UCC operator in the set of gates available to most quantum computers, one generally has to perform a Trotter decomposition\cite{Trotter1959, Suzuki1976} of the cluster operator,
\begin{equation}
 e^{\hat \tau} = \lim_{\rho \rightarrow \infty}\Big(\prod_i e^{\hat\tau_i/\rho}\Big)^\rho.
\end{equation}
Usually, rather than working in the large $\rho$ limit, $\rho$ is taken to be unity, leading to a disentangled from of the UCC Ansatz, which has been found to be as expressive as the original algorithm under most circumstances.\cite{Evangelista2019}

Because of this decomposition, UCC-based circuits often tend to be prohibitively deep for NISQ devices. Various alternative algorithms that attempt to balance circuit depth against potential energy surface complexity have been defined, including symmetry preserving Ans\"atze,\cite{Gard2020} adaptive Ans\"atze\cite{Grimsley2019b} and methods like qubit coupled cluster (QCC)\cite{Ryabinkin2018,Ryabinkin2020} and Givens-rotation-based Ans\"atze\cite{Arrazola2022} which avoid some of the complexities of fermionic excitation operators.
\subsection{Projective Quantum Eigensolver}
As an alternative to variational algorithms, one can solve electronic structure problems using projective approaches. In classical settings, this technique is most commonly employed in Projector Monte Carlo (PMC) algorithms, which will be discussed in more detail in \cref{sec:qmc}.

Consider a trial wavefunction expressed as \mbox{$\ket{\Psi_0} =\hat U (\boldsymbol\theta)\ket{\phi_0}$}, where $\ket{\phi_0}$ is a simple reference wavefunction such as a Hartree--Fock determinant and $\hat U(\mathbf{\boldsymbol\theta})$ is some parametrised unitary operator.
If this were an exact solution to the Schr\"odinger equation, then
\begin{equation}
\hat H \hat U \ket{\phi_0} = E \hat U\ket{\phi_0}
\end{equation}
and
\begin{equation}
\hat U^\dagger \hat H \hat U \ket{\phi_0} = E \ket{\phi_0}.
\label{eq:linked}
\end{equation}
Projecting \Cref{eq:linked} onto the reference wavefunction $\ket{\phi_0}$ gives
\begin{equation}
    \braket{\phi_0|\hat U^\dagger \hat H \hat U |\phi_0} = E_\mathrm{PQE},
    \label{eq:epqe}
\end{equation}
while projecting onto the complete set of $\ket{\phi_\mathbf{i}}$ orthogonal to $\ket{\phi_0}$ gives a set of equations for the residuals
\begin{equation}
 r_\mathbf{i} = \braket{\phi_\mathbf{i}\hat| U^\dagger \hat H \hat U |\phi_0},
\end{equation}
which all satisfy
\begin{equation}
    r_\mathbf{i} = 0.
     \label{eq:resid}
\end{equation}
If $\hat U\ket{\phi_0}$ is the true ground state wavefunction, \Cref{eq:resid} holds for all states in the Hilbert space. In general, it can be enforced only for a subset of states $\ket{\phi_\mathbf{i}}$, corresponding to parameters in $\hat U(\boldsymbol \theta)$.  Solving \Cref{eq:epqe,eq:resid} for this subset leads to a Projective Quantum Eigensolver (PQE)\cite{Stair2021} algorithm.

If $\hat U$ is a UCC cluster operator and $\ket{\phi_0}$ is the Hartree--Fock determinant, \Cref{eq:epqe,eq:resid} can then be solved using a quasi-Newton iterative scheme such as\cite{Stair2021} 
\begin{equation}
\theta_\mathbf{i}^{(n+1)} =\theta_\mathbf{i}^{(n)} + \frac{r_\mathbf{i}}{\Delta_\mathbf{i}}
\end{equation}
with ${\Delta_\mathbf{i}} = \epsilon_i + \epsilon_j + ... - \epsilon_a - \epsilon_b - ...$, where $\epsilon_i$ are Hartree--Fock orbital energies and $\ket{\phi_{\mathbf i}}$ labels $\ket{\phi_{ij...}^{ab...}}$. This algorithm is appealing for NISQ device use because, unlike conventional VQE which requires gradient computations, it only requires the residuals $r_\mathbf{i}$ which can be computed as expectation values at the same cost as the overall energy, since
\begin{align}
\begin{split}
r_\mathbf{i} = &\braket{\Omega_\mathbf{i}(\pi/4)|\hat U^\dagger \hat H \hat U|\Omega_\mathbf{i}(\pi/4)} \\
&- \frac{1}{2}\braket{\phi_0|\hat U^\dagger \hat H \hat|\hat U \phi_0} -\frac{1}{2}\braket{\phi_\mathbf{i}|\hat U^\dagger \hat H \hat U|\phi_\mathbf{i}}
\end{split}
\label{eq:exp_resid}
\end{align}
where $\ket{\Omega_\mathbf{i}(\theta)} = e^{\theta \hat \tau_\mathbf{i}}\ket{\phi_0} = e^{\theta (\hat a_\mathbf{i} - \hat a_\mathbf{i}^\dagger)}\ket{\phi_0} = {\cos{\theta}\ket{\phi_0} + \sin{\theta}\ket{\phi_\mathbf{i}}}$.

\subsection{Quantum Monte Carlo}\label{sec:qmc}

As mentioned before, the main use of projective methods in conventional quantum chemistry is in PMC algorithms, which are based on the imaginary-time Schr\"odinger equation,
\begin{equation}
    \frac{\partial \ket{\Psi}}{\partial \beta} = - \hat H \ket{\Psi},
    \label{eq:im_sch}
\end{equation}
where $\beta$ is imaginary time and $\hat H$ is the Hamiltonian of the system of interest. The lowest-energy solution to \Cref{eq:im_sch}, $\ket{\Psi_0}$, is given by
\begin{equation}
    \ket{\Psi_0} = \lim_{\beta \rightarrow \infty} e^{-\beta (\hat H - E_0)}\ket{\Psi(0)},
\end{equation}
where $\ket{\Psi(0)}$ is some initial trial wavefunction such that $\braket{\Psi_0|\Psi(0)}\neq0$ and $E_0$ is the lowest eigenvalue of $\hat H$. The time-evolution operator can also undergo a Trotter expansion, giving
\begin{equation}
    \ket{\Psi_0} = \lim_{n \rightarrow \infty} (e^{-\Delta \beta (\hat H - E_0)})^n \ket{\Psi(0)},
\end{equation}
where $\Delta \beta = \beta/n$. For small $\Delta \beta$, the exponential can be approximated by a linear expansion, leading to the form of the imaginary time propagator most commonly used in projector methods.
\begin{equation}
    \ket{\Psi_0} = \lim_{n \rightarrow \infty} (1 - \Delta\beta (\hat H - E_0))^n \ket{\Psi(0)}
    \label{eq:proj}
\end{equation}
Considering a single time-step $\Delta \beta$, one can write
\begin{equation}
        \ket{\Psi(\beta + \Delta \beta} = (1 - \Delta\beta (\hat H - E_0)) \ket{\Psi(\beta)},
\end{equation}
which describes the step-wise imaginary time evolution of the trial wavefunction $\ket{\Psi}$. As in PQE, this equation can be projected onto the various Slater determinants in the Hilbert space to give
\begin{equation}
    \braket{\phi_\mathbf{i}|\Psi(\beta + \Delta \beta)} = \braket{\phi_\mathbf{i}|1 - \Delta\beta (\hat H - E_0)|\Psi(\beta)}
    \label{eq:pmc}
\end{equation}
Provided $\ket{\Psi}$ is parametrised with respect to $\{\boldsymbol \theta\}$ such that $\braket{\phi_\mathbf{i}|\Psi(\boldsymbol{\theta};\beta)} = \theta_\mathbf{i} + \mathcal{O}(\theta^{1 + \epsilon})$ for any $\epsilon > 0$,  \Cref{eq:pmc} can be approximated as
\begin{equation}
    \theta_\mathbf{i}(\beta + \Delta \beta) = \theta_\mathbf{i}(\beta) - \Delta\beta\braket{\phi_\mathbf{i}|\hat H-E_0|\Psi(\beta)}
\end{equation}
or
\begin{equation}
  \theta_\mathbf{i}(\beta + \Delta \beta) = \theta_\mathbf{i}(\beta) - \Delta\beta\Big(c_\mathbf{i} (H_\mathbf{ii} - E_0) + \sum_{\mathbf{j} \neq \mathbf{i}} c_\mathbf{j}H_\mathbf{ij}\Big),  
  \label{eq:pop_dyn}
\end{equation}
where $c_\mathbf{i}$ is the configuration interaction (CI) coefficient of $\ket{\phi_\mathbf{i}}$ in $\ket{\Psi}$ and $H_\mathbf{ij} = \braket{\phi_\mathbf{i}|\hat H|\phi_\mathbf{j}}$. We note that this is equivalent to 
\begin{equation}
\theta_\mathbf{i}(\beta + \Delta \beta) = \theta_\mathbf{i}(\beta) - \Delta\beta r_\mathbf{i},
\label{eq:param_resid}
\end{equation}
where the residual is now defined as $r_\mathbf{i} = \braket{\phi_\mathbf{i}|\hat H-E_0|\Psi(\beta)}$. The energy appearing in the residual is due to the effective use of an unlinked form of the cluster equations.

In Quantum Monte Carlo (QMC) methods, \Cref{eq:pop_dyn} is interpreted as governing the population dynamics of a set of particles (``walkers", ``psips" or ``excips") living in the Hilbert space. Rather than iterating these equations exactly, their terms are sampled stochastically and can be described by three processes:\cite{Booth2009}
\begin{itemize}
\item \textbf{Spawning} of particles from determinant $\ket{\phi_\mathbf{j}}$ to parameter $\theta_\mathbf{i}$ corresponds to the off-diagonal action of the Hamiltonian.
\begin{equation} 
\delta \theta_\mathbf{i}^{(s, k)} = - \Delta\beta \frac{c_\mathbf{j}H_\mathbf{ij}}{p_\mathrm{select}(\phi_\mathbf{j}) p_\mathrm{gen}(\mathbf{i}|\mathbf{j}) n_\mathrm{attempts}},
\label{eq:pspawn}
\end{equation}
where $p_\mathrm{select}(\phi_\mathbf{j})$ is the probability of selecting $\phi_\mathbf{j}$ as the source of the spawn, $p_\mathrm{gen}(\mathbf{i}|\mathbf{j})$ is the conditional probability of selecting $\theta_\mathbf{i}$ as the target of a spawn given it originates from $\phi_\mathbf{j}$ and $n_\mathrm{attempts}$ is the number of spawning attempts. $k$ is the index of the current attempt and $s$ denotes that this is a spawning event.
\item \textbf{Death} or cloning from $\ket{\phi_\mathbf{i}}$ to parameter $\theta_\mathbf{i}$ corresponds to the diagonal action of the Hamiltonian.
\begin{equation} 
\delta \theta_\mathbf{i}^{(d, k)} = - \Delta\beta \frac{c_\mathbf{i} (H_\mathbf{ii} - S)}{p_\mathrm{select}(\phi_\mathbf{i}) n_\mathrm{attempts}},
\end{equation}
where $d$ denotes that this is now a death event.
\item \textbf{Annihilation} ensures that the two equivalent oppositely signed versions of the wavefunction do not proliferate simultaneously and corresponds to collecting all contributions to $\theta_\mathbf{i}$ into a single value.
\begin{equation}
\theta_\mathbf{i} = \theta_\mathbf{i} + \sum_k^{n_\mathrm{attempts}}\delta \theta_\mathbf{i}^{(s, k)} + \sum_k^{n_\mathrm{attempts}}\delta \theta_\mathbf{i}^{(d, k)}
\end{equation}
\end{itemize}


\begin{figure*}
    \centering
\subfloat[][ ]{\begin{quantikz}
\lstick{$\ket{0}$} & \qw& \gate{H}& \ctrl{1}  &\octrl{1}&\gate{H}&\qw\\
\lstick{$\ket{0}$}&\gate{\prod_{i \in \mathrm{occ. 0}} X_i}\qwbundle[alternate]{}&\qwbundle[alternate]{} &\gate{U}\qwbundle[alternate]{}&\gate{\prod_{i \in \mathrm{act. orb}} X_i} \qwbundle[alternate]{}&\qwbundle[alternate]{}&\qwbundle[alternate]{}
\end{quantikz}}\\
\subfloat[][ ]{\begin{quantikz}
\lstick{$\ket{0}$} &  \gate{H}& \octrl{1}  &\ctrl{1}&\qw&\ctrl{1}&\qw&\ctrl{1}&\gate{H}&\qw\\
\lstick{$\ket{0}$}&\qwbundle[alternate]{} &\gate{\prod_{i \in \mathrm{occ. 0}} X_i}\qwbundle[alternate]{}&\gate{\prod_{i \in \mathrm{occ, 1}} X_i} \qwbundle[alternate]{}&\qwbundle[alternate]{}&\gate[2]{\mathrm{SWAP}}\qwbundle[alternate]{}&\gate{U}\qwbundle[alternate]{}&\gate[2]{\mathrm{SWAP}}\qwbundle[alternate]{}&\qwbundle[alternate]{}&\qwbundle[alternate]{}\\
\lstick{$\ket{0}$}&\qwbundle[alternate]{}&\qwbundle[alternate]{}&\qwbundle[alternate]{}&\qwbundle[alternate]{}&\qwbundle[alternate]{}&\qwbundle[alternate]{}&\qwbundle[alternate]{}&\qwbundle[alternate]{}&\qwbundle[alternate]{}
\end{quantikz}}
    \caption{\small \justifying Quantum circuits for measuring the overlap $s_\mathrm{i} = \braket{\phi_\mathbf{i}|\hat U \phi_0}$. The first is a modified Hadamard test,\cite{quantumalgorithms} which uses a single ancilla qubit and a controlled version of the unitary $\hat U$. The second was proposed by Huggins \textit{et al.}\cite{Huggins2020} and needs $(n_\mathrm{qubits} + 1)$ ancilla but no controlled $\hat U$.}
    \label{fig:circs}
\end{figure*}

We note that in the death step the unknown true ground state energy $E_0$ has been replaced by the shift $S$. This acts as a population control parameter and in standard QMC procedures eventually converges to a stochastic estimator of $E_0$.\cite{Booth2009} The projected energy
\begin{equation}
    E_\mathrm{proj} = \frac{\braket{\phi_0|\hat H| \Psi(\beta)}}{\braket{\phi_0|\Psi_(\beta)}}.
    \label{eq:e_proj}
\end{equation}
is normally used as a second way to estimate $E_0$.

In the most popular QMC algorithm, Full Configuration Interaction QMC (FCIQMC),\cite{Booth2009} the wavefunction is parametrised directly as a CI expansion so $\theta_\mathbf{i} = c_\mathbf{i}$, which serves perhaps to justify the names given to the steps in the Monte Carlo propagation algorithm. Alternative formulations have been developed using Coupled Cluster (CC)\cite{Thom2010} and UCC\cite{Filip2020} parametrisations of the wavefunction. These methods use very similar propagators, however  one only has access to the cluster amplitudes $t_\mathbf{i}$. Therefore the corresponding CI coefficients must also be sampled stochastically, leading to additional complexity in the selection algorithm, which is responsible for a significant portion of the computational cost of these algorithms. 

In FCIQMC it is easy to cycle through all determinants present in the wavefunction or sample them uniformly. In contrast, in Coupled Cluster Monte Carlo (CCMC) a doubly excited determinant $\ket{\phi_\mathbf{i}} = \ket{\phi_{ij}^{ab}}$ will have a CI coefficient given by
\begin{equation}
    c_{ij}^{ab} = t_{ij}^{ab} + t_{i}^at_j^b - t_i^bt_j^a,
\end{equation}
where $t_{i}^a$, $t_{ij}^{ab}$ are the amplitudes of the corresponding single and double excitation operators $\hat a_i^a$ and $\hat a_{ij}^{ab}$ in the cluster operator $\hat T$. Therefore, when selecting a determinant for spawning or death, one must consider both composite (e.g. $t_{i}^at_j^b$) and non-composite ($t_{ij}^{ab}$) contributions (``clusters") to it. The standard way to do this in CCMC\cite{Thom2010, Scott2017} is to:
\begin{enumerate}
    \item Select a cluster size $s$ with some probability $p_s$;
    \item (Optionally) select a particular type of cluster of size $s$ by imposing constraints on the excitation levels of the individual excitors involved;\cite{Scott2017}
    \item Select $s$ excitors to form the cluster, each with probability $p_e \propto t_e$;
    \item ``Collapse" the excitors to determine the corresponding determinant (e.g. $\hat a_i^a \hat a_j^b \ket{\phi_0} = \pm\ket{\phi_{ij}^{ab}}$).
\end{enumerate}

A variety of selection algorithms have been developed for CCMC  to better importance sample the underlying wavefunction.\cite{Scott2017} Nevertheless, this remains one of the main challenges in CCMC, particularly when expanded to high truncation levels\cite{Neufeld2017} or multi-reference formulations.\cite{Filip2019} 

To further complicate matters, in unitary CCMC (UCCMC)\cite{Filip2020} the maximum cluster size is formally infinite, due to the alternation of excitation and de-excitation operators in the clusters. While results converge relatively quickly with maximum cluster size, the inclusion of large composite clusters leads to instabilities in the Monte Carlo propagation.


\subsection{Folded-spectrum} 
Methods like VQE, PQE and standard PMC are all intended to find the ground state of a system of interest. Some excited states can be found with these techniques by enforcing orthogonality between the state of interest and the ground state, for example by considering different symmetry sectors of the Hilbert space.\cite{Nakanishi2019} However, obtaining general excited states is not-trivial.

One approach to directly target excited states is by using the folded-spectrum approach,\cite{Peruzzo2014,Cao2019,CadiTazi2023} in which the system Hamiltonian $\hat H$ is replaced by
\begin{equation}
\hat H_\mathrm{fs} (\omega) =  (\hat H - \omega)^2,
\label{eq:fs_ham}
\end{equation}
which has the same eigenstates as the original $\hat H$. However, the eigenvalues become $(E_i - \omega)^2$ so the ground state of $\hat H_\mathrm{fs}(\omega)$ becomes the eigenstate with true energy $E_i$ closest to $\omega$. Therefore, by changing $\omega$ one can target any stationary state of the system with conventional variational or projective techniques. Similar approaches have been used in the conventional Variational Monte Carlo (VMC) community under the name of state-specific variance optimisation.\cite{Umrigar1988, Zhao2016,Shea2017,Cuzzocrea2020, Otis2023}

In variational folded-spectrum approaches, rather than the energy one minimises the quantity
\begin{equation}
\braket{E}_\mathrm{fs} (\omega) = \frac{\braket{\Psi(\boldsymbol{\theta})|(\hat H - \omega)^2| \Psi(\boldsymbol{\theta})}}{\braket{\Psi(\boldsymbol{\theta})|\Psi(\boldsymbol{\theta})}},
\end{equation}
while in projective approaches one updates the parameters using the residuals
\begin{equation}
r_\mathbf{i} = \braket{\phi_\mathbf{i}|(\hat H - \omega)^2 -E_\mathrm{fs,0}| \Psi(\boldsymbol{\theta})}.
\label{eq:res_fs}
\end{equation}
While the algorithms employed are equivalent to the standard ground-state approaches, the folded-spectrum Hamiltonian is more complex than the original Hamiltonian. For a standard chemistry Hamiltonian
\begin{equation}
    \hat H = \sum_{pq} f_{pq} \hat p^\dagger \hat q + \sum_{pqrs}g_{pqrs}\hat p^\dagger \hat q ^\dagger \hat r \hat s,
    \label{eq:ferm_hamil}
\end{equation}
the corresponding folded-spectrum Hamiltonian will include up to four-body terms, 
\begin{equation}
    \sum_{pqrstuvw}g_{pqrs}g_{tuvw}\hat p^\dagger \hat q ^\dagger \hat t^\dagger \hat s^\dagger \hat r \hat s \hat v \hat w,
\end{equation}
which will be significantly more expensive to compute than the one- and two-body terms of the original Hamiltonian.
\begin{figure*}
\centering
\begin{quantikz}
\lstick{$|a\rangle$} & \qw & \qw & \qw  &\ctrl{4} & \qw & \qw &\qw & \qw \\
\lstick{$|q_0\rangle$} & \qw & \qw & \qw  &\qw & \qw & \qw  & \qw& \qw \\
\lstick{$|q_1\rangle$} & \gate{H} & \ctrl{1}  & \qw & \qw & \qw & \ctrl{1}   & \gate{H} & \qw \\
\lstick{$|q_2\rangle$} & \qw & \targ{} & \ctrl{1}   & \qw & \ctrl{1} & \targ{} & \qw  & \qw \\
\lstick{$|q_3\rangle$} & \gate{R_x(\frac{3\pi}{2})} &\qw & \targ{}  & \gate{R_z(\theta)} & \targ{} & \qw & \gate{R_x(\frac{\pi}{2})} & \qw \\
\end{quantikz}
\caption{\small Controlled Pauli gadget encoding $e^{-i\frac{\theta}{2}( X_1 \otimes Z_2\otimes Y_3)}$.}
\label{fig:controlled_pauli}
\end{figure*}
\section{Quantum Algorithms}
\subsection{Unlinked Residual Measurement Implementation}\label{sec:algorithm}
The QMC methods discussed in Sec. \ref{sec:qmc} are all based on computing the unlinked residual 
\begin{equation}
    r_\mathbf{i} = \braket{\phi_\mathbf{i}|\hat H-S|\Psi(\beta)}
\end{equation}
and the overlap
\begin{equation}
    s_\mathrm{i} = \braket{\phi_\mathbf{i}|\Psi(\beta)},
\end{equation}
As detailed in Sec. \ref{sec:qmc}, exponential Ansatz-based QMC algorithms require complex selection algorithms to estimate $r_\mathbf{i}$, as computing it directly would require storing the full CI wavefunction, which is not memory-efficient. For large cluster sizes, this often becomes the limiting factor on the speed and rate of convergence of a calculation. Therefore, using a quantum device to compute these terms exactly without the memory overhead would be a significant boon to these methods. 

The standard implementation of the UCC Ansatz in the quantum circuit representation is to map Fermionic creation and annihilation operators onto strings of Pauli $X$, $Y$ and $Z$ matrices (or gates). There are multiple standard mapping schemes,\cite{Jordan1928,Bravyi2002,Seeley2012,Tranter2015} of which we employ here the Jordan--Wigner approach,\cite{Jordan1928} in which

\begin{equation}
\begin{split}
\hat a_p^\dagger = \sigma_p^+ \bigotimes_{q>p} Z_q \\
\hat a_p = \sigma_p^- \bigotimes_{q>p} Z_q,
\end{split}
\end{equation}
where the product of $Z$ matrices encodes the required parity to preserve fermionic anti-commutation relations. Using this mapping leads to the following forms for single and double excitation operators:

\begin{equation}
\hat \tau_{i}^{a} = t_i^{a}(a^\dagger_i a_a - a^\dagger_a a_i) = \frac{it_{i}^{a}}{2} \bigotimes_{k=i+1}^{a-1} Z_k(Y_iX_a - X_iY_a),
\end{equation}
\begin{align}
\begin{split}
\hat \tau_{ij}^{ab} &= t_{ij}^{ab}(a^\dagger_a a^\dagger_b a_i a_j - a^\dagger_j a^\dagger_i a_a a_b)  \\
 &=\frac{it_{ij}^{ab}}{8}  \bigotimes_{k=i+1}^{j-1} Z_k  \bigotimes_{l=a+1}^{b-1} Z_l (X_i X_j Y_a X_b + Y_i X_j Y_a Y_b \\
 &+X_i Y_j Y_a Y_b + X_i X_j X_a Y_b  - Y_i X_j X_a X_b \\
 &- X_i Y_j X_a X_b - Y_i Y_j Y_a X_b - Y_i Y_j X_a Y_b).
\end{split}
\end{align}
The exponential of each of these strings of Pauli matrices can be computed using a set of gates known as a Pauli gadget,\cite{Cowtan2020} a controlled version of which is shown in \Cref{fig:controlled_pauli}.

It has previously been shown that overlap with complex wavefunctions can be efficiently calculated using a quantum computer\cite{Huggins2022} in the context of auxilliary-field QMC (AFQMC)\cite{Zhang2003} trial wavefunctions. As there is no immediate expectation-value form similar to \Cref{eq:exp_resid} for $r_\mathbf{i}$ and $s_\mathbf{i}$ computing such quantities is possible, but requires auxiliary qubits. Two alternative circuit constructions to obtain $r_\mathbf{i}$ and $s_\mathbf{i}$ are given in \Cref{fig:circs}. 

The first uses fewer ancilla qubits, but it also requires a controlled version of the unitary $\hat U$, which decomposes into a larger number of two-qubit gates than its uncontrolled couterpart, regardless of the nature of $\hat U$. As a significant fraction of NISQ device error is due to two-qubit gates, the second circuit may be preferable on hardware. However, we find that when working with simulators the $2^{n_\mathrm{qubits}}$-fold increase of the size of the qubit Hilbert space leads to worse performance, so work presented here employs the circuit in \Cref{fig:circs} (a). 
 Additionally, in the particular case of the Pauli gadget-based unitary coupled cluster Ansatz, controlled-$\hat U$ only requires one additional controlled rotation gate per Pauli gate string, making its implementation relatively painless.

Labelling the overall state in the qubit register of the circuit in \Cref{fig:circs} (a) as $\ket{\Phi}$, the expectation value of $Z$ on the ancilla is given by
\begin{widetext}
\begin{equation}
        \braket{\Phi|\hat Z_\mathrm{anc}|\Phi} =\frac{1}{4} (\braket{\phi_i| \hat U\phi_0} + \braket{\phi_0 \hat U^\dagger|\phi_i} + \braket{\phi_i| \hat U \phi_0} + \braket{\phi_0 \hat U^\dagger|\phi_i}
        = \frac{1}{4} (4 \times \Re{\braket{\phi_i|\hat U\phi_0}}) = \Re{\braket{\phi_i|\hat U\phi_0}},
\end{equation}
\end{widetext}
so $\braket{\Phi|\hat Z_\mathrm{anc}|\Phi} = \Re{(s_\mathrm{i})}$. The $Y$-expectation value similarly gives the imaginary part of the overlap. Additionally, applying the Hamiltonian to the system register and measuring the expectation value gives the corresponding parts of the residual,
\begin{equation}
\braket{\Phi|\hat Z_\mathrm{aux}\hat H|\Phi} = \Re{\braket{\phi_i|\hat H|\hat U\phi_0}} = \Re{(r_\mathbf{i})}.
\end{equation}
Residuals computed in this way can be employed in a projector method based on \Cref{eq:param_resid}, leading to an unlinked PQE algorithm. Given noiseless, exact values for $r_\mathbf{i}$ and setting the shift to track the instantaneous projected energy, such an algorithm would converge directly to the ground state, without need for further Monte Carlo estimation. However, as data is obtained from a realistic quantum device in finite samples, it is intrinsically noisy even when ignoring NISQ-era device noise. Therefore, it is perhaps wiser to consider the measured $r_\mathbf{i}$ as stochastic estimates of the underlying quantities and therefore use them in a QMC-like algorithm as outlined below:
\begin{enumerate}
    \item On a classical machine, represent the wavefunction $\hat U (\boldsymbol{\theta})\ket{\phi_0}$ as a distribution of walkers with populations $N_0$ corresponding to $\ket{\phi_0}$ and $N_i = \theta_i N_0$ corresponding to each parameter.
    \item Encode the wavefunction $\hat U (\boldsymbol{\theta})\ket{\phi_0}$ on a quantum device and measure the residuals $r_\mathrm{i}$
    \item Update each $N_\mathbf{i} = N_\mathbf{i} - N_0 \times \delta\beta \times r_\mathbf{i}$
    \item Store the projected energy $E_\mathrm{proj} = r_0 + S * s_0$.
    \item Once a threshold population $N_\mathrm{tot} = \sum_{\mathbf{i}=0}^{N_\mathrm{param} +1} N_\mathbf{i}$ has been reached, allow the shift to vary to maintain this population constant. In conventional QMC this is done by\cite{Booth2009}
    \begin{equation}
        S(\beta) = S(\beta - A\delta\beta) - \frac{\zeta}{A\delta\beta}\mathrm{ln}\frac{N_\mathrm{tot}(\beta)}{N_\mathrm{tot}(\beta - A\delta\beta)}.
        \label{eq:shift}
    \end{equation}
    \item Repeat steps 2-5 until sufficiently many samples of $E_0$ and $S$ have been accumulated to obtain accurate Monte Carlo estimates.
\end{enumerate}
On the one hand, as discussed in Sec. \ref{sec:res}, this approach has a series of benefits relative to a ``deterministic" PQE approach, as it is resilient to noise in the quantum measurement, allowing the use of low shot numbers to obtain accurate results. Additionally, the shift provides a second independent estimator for the energy which can be used to more clearly identify when convergence has been reached. On the other hand, it is more efficient than a fully classical Quantum Monte Carlo algorithm as it both removes the need for cluster sampling and merges the death and spawning steps into a single process.

One of the main benefits of QMC methods over conventional algorithms is that they take advantage of sparsity in the Hamiltonian to store a compressed version of the wavefunction, thereby decreasing classical memory overheads. In the following section we explore other QMC-inspired ideas that can be used to further decrease the memory or quantum resource costs of the algorithm detailed above.

\subsection{Stochastic additions}

\subsubsection{Sampling the Hamiltonian}

The second-quantised Hamiltonian in \Cref{eq:ferm_hamil} can be mapped onto a corresponding qubit Hamiltonian using the same techniques discussed for excitation operators. This leads to an operator of the form
\begin{equation}
    \hat H = \sum_k h_k \hat P_k,
\end{equation}
where $\hat P_k$ are strings of Pauli operators. When computing the residuals $r_\mathbf{i}$ one must in principle measure each of the terms of the Hamiltonian separately, as they require different post-rotations to be applied to the circuit before measurement. They can be further grouped\cite{McClean2016, Kandala2017} as 
\begin{equation}
    \hat H = \sum_k \big( \sum_{j \in G_k}  h_{kj} \hat P_{kj}\big), 
\end{equation}
where $[\hat P_{ki}, \hat P_{lj}] \propto \delta_{kl}$. Terms in the same group $G_k$ commute, so they may be measured simultaneously using the same post-rotations, but this is nevertheless the bulk of the quantum computation required. Additionally, optimal grouping of Pauli terms is in general a hard problem.\cite{Gokhale2020} For the two-body quantum chemistry Hamiltonian there is a simple heuristic partitioning, based of the qubitwise commutativity of Pauli strings. This leads to one large group comprising the identity string and all strings containing only $Z$ gates and many small groups of strings containing some $X$ and $Y$ gates.

One could employ an approach similar to that used in classical QMC algorithms and only select a subset of the groups to measure at each time-step. This decreases the number of circuit measurements needed per time-step, but introduces further noise in the estimate of $r_\mathbf{i}$. However, provided no systematic bias is caused, this noise will be averaged out over the course of the imaginary-time propagation.

There exist various possible Hamiltonian term selection schemes, the simplest of which would be to sample uniformly from the set of Pauli terms or groups. Many of these terms will have negligible weights however, so in order to better capture the action of the Hamiltonian in a small number of samples, some form of importance sampling is recommended. One straightforward approach is to select each Pauli group with probability proportional the group weight $W_k = \sum_{j \in G_k} |h_{kj}|$,

\begin{equation}
    p(k) = \frac{W_k}{\sum_k W_k}
\end{equation}

For most quantum chemistry Hamiltonians investigated in this paper, we find that the weight of the group containing the empty Pauli string and all $Z$-only strings, which we will label $G_0$, strongly dominates the resulting probability distribution. Therefore, under this selection scheme, it is likely that for small numbers of samples it will be \textit{over-} rather than \textit{under}sampled. We also note that this term is qualitatively different from the other groups, as it describes the diagonal action of the Hamiltonian, while the others all correspond to off-diagonal contributions. Therefore, to obtain a more balanced description, we propose a scheme in which the diagonal group $G_0$ is always evaluated, while the other groups are selected with probability
\begin{equation}
    p'(k) = \frac{W_k}{\sum_{k\neq 0} W_k}
\end{equation}
\subsubsection{Sampling the wavefunction}

The natural candidate Ansatz for this type of quantum-enhanced QMC algorithm is unitary coupled cluster. One of the main challenges with using this Ansatz on NISQ devices is that it leads to deep circuits, with execution times that easily exceed the coherence times for current machines. We have previously shown that by using insight from a fully classical UCCMC calculation,\cite{Filip2020} one can significantly shorten circuits while preserving accuracy by removing low-weight excitors from the Ansatz.\cite{Filip2022}

This circuit simplification can be done on-the-fly during the propagation, by stochastically rounding the parameter values before encoding the wavefunction,
\begin{equation}
    \widetilde \theta_\mathbf{i} = \begin{cases}
    \theta_\mathbf{i},\ \mathrm{if}\ \frac{\sum_\mathbf{j}^\mathbf{i}\theta_\mathbf{j}}{\sum_\mathbf{j}^{\mathbf{i}_\mathrm{max}}\theta_\mathbf{j}} \leq p\\
    0,\ \mathrm{otherwise},
    \end{cases}
\end{equation}
where $p$ is a uniform random sample from the interval $[0,1)$ and the parameters $\theta_\mathbf{i}$ are ordered by decreasing amplitude. Once again, this is expected to increase fluctuations in the estimators at each time-step. The issue of potential systematic biases is explored further in Sec. \ref{sec:res}.

\subsubsection{Spawning}\label{sec:spawn}
In the methodology detailed above, all parameters are updated at every time-step. However, one can take the approach commonly employed in QMC algorithms and only update a stochastically selected subset of the populations. In particular, for spawning in conventional QMC algorithms one can use a variety of excitation generation algorithms.\cite{Holmes2016,Neufeld2019} 

In principle, one could spawn onto connected determinants uniformly. However the spawning probability is proportional to $\frac{H_\mathbf{ij}}{p_\mathrm{gen}(\mathbf{i|j}) }$ (see \Cref{eq:pspawn}), so for appropriate importance sampling of this process it is desirable that ${p_\mathrm{gen}(\mathbf{i|j}) } \propto {H_\mathbf{ij}}$. Different algorithms have been designed to approximately achieve this, using heat-bath sampling\cite{Holmes2016} or mathematical bounds on the value of ${H_\mathbf{ij}}$.\cite{Neufeld2019}

With the residuals computed according to the quantum algorithm in Sec. \ref{sec:algorithm}, death and spawning are no longer separate processes, so the parameters to be updated should be generated with 
\begin{equation}
{p_\mathrm{gen}(\mathbf i) } \propto \braket{\phi_\mathbf{i}|\hat H - S |\Psi}.
\end{equation}

Generating samples from $\hat H - S \ket{\Psi}$ is non-trivial even with access to a quantum computer and a qubit representation of the wavefunction and Hamiltonian. In particular, the circuits described in \Cref{fig:circs}, which can be used to estimate the overlap itself, cannot be used for this purpose. However, one can efficiently generate determinants with probability  $\propto |\braket{\phi_\mathbf{i}|\hat H - S|\Psi}|^2$ by preparing the state $\ket{\Psi}$, applying one of the operators $\hat P_k$ to it with probability $\frac{|h_k|}{\sum_k|h_k|}$ and then measuring the qubit register. While this is not a perfect importance sampling of the Hamiltonian coupling, it will match its minima and maxima and should therefore represent an improvement over uniform sampling. These circuits do not require ancilla or additional controlled operations, so they are less costly to implement than those needed to compute the residuals. The number of shots needed is equal to the number of determinants one wants to generate, so they would add little computational overhead.

While stochastic parameter updates may be valuable for larger systems, for the problems considered in this paper it is easy to store and cycle through all the relevant parameters and therefore probabilistic spawning is not employed.

\subsection{Folded spectrum considerations}
All the methods described above can equally be applied to the folded spectrum Hamiltonian, by replacing the residual with that in \Cref{eq:res_fs}. As the squared Hamiltonian operator contains many-body terms of higher order than the original Hamiltonian, this will lead to more terms in the Pauli string representation of the corresponding qubit operator as well. However, as some of these will commute with the terms in the original operator, the number of different groups will see a smaller increase. 
Nevertheless, the possibility of Hamiltonian sampling is particularly promising in this case to mitigate the additional costs.

\begin{figure}
    \centering
    \includegraphics[width=0.5 \textwidth, trim = 0 1cm 0 1.5cm, clip]{H3+-S-2.0.pdf}\\
    \includegraphics[width=0.5 \textwidth, trim = 0 1cm 0 2cm, clip]{H3+-E-2.0.pdf}
    \caption{\small \justifying Shift (top) and projected energy (bottom) as a function of imaginary time, obtained from $10^1 - 10^4$ samples of the residual circuits for each Pauli term in the Hamiltonian,  for H$_3^+$ at $r = 2.0$ \r{A}. Noise decreases as the number of shots increases.}
    \label{fig:shot_noise}
\end{figure}

As molecules are pulled apart electronic states approach one another and often become degenerate in the fully dissociated limit. Therefore, if one is searching for two different states which are close in energy (which will remain so after the folding procedure), noise in the wavefunction representation becomes potentially problematic as it may be sufficient to induce a switch from one state to another. This is observed in conventional CCMC, where large fluctuations will sometimes cause a calculation to collapse onto a different, often unphysical state.\cite{Scott2019,Scott2020} Therefore, the range of stochastic techniques that can be used in this scenario may be more limited than in ground state calculations.

\section{Results}\label{sec:res}
All results in this paper are obtained from state-vector or shot-based quantum circuit simulation using the $\mathrm{t}\ket{\mathrm{ket}}$ platform.\cite{Sivarajah2021} Necessary atomic orbital Hamiltonian integrals were obtained from PySCF\cite{Sun2020} and FCI calculations were performed using HANDE-QMC.\cite{Spencer2019} All calculations used the UCCSD Ansatz. The main systems considered are the H$_2$, H$_3^+$, H$_3$ and H$_4$ linear molecules in the STO-3G basis set,\cite{Hehre1969} which serve as a useful proof of concept for the methods described here. These molecules are small enough that the corresponding Fock space vectors can be easily manipulated on a classical computer, while covering a range of chemical behaviour. Shot-based simulations are obtained by taking $N_\mathrm{shots}$ samples from the probability distribution corresponding to the wavefunction encoded in a given quantum circuit. 
\begin{table}
    \centering
    \begin{tabular}{c|c|c|c|c|c}
       $N_\mathrm{shots}$  &  $\braket{E}/E_\mathrm{h}$ & $\sigma (E)/E_\mathrm{h}$ & $\braket{S}/E_\mathrm{h}$ & $\sigma(S)/E_\mathrm{h}$&$E_\mathrm{FCI}/E_\mathrm{h}$\\
       \hline
         10 & -0.111 & 0.009 & - & - &\multirow{4}{*}{-0.11586}\\
         100 & -0.116& 0.003& -0.114& 0.004& \\
         1000 & -0.116.7& 0.0009 & -0.117& 0.001& \\
         10000 & -0.1163 & 0.0003& -0.1155 & 0.0004& \\
         \hline
    \end{tabular}
    \caption{\small \justifying Mean and standard deviation of the shift and projected energy as the number of shots is increased, for H$_3^+$ at $r = 2.0$ \r{A}}
    \label{tab:variance}
\end{table}
In order to reduce computation time required for larger systems, shot-based simulation is at times replaced by state-vector simulation with added Gaussian noise of mean $\mu = 0$ and standard deviation $\sigma$. The average behaviour of such simulations should be very similar to true shot-based results for large enough $N_\mathrm{shots}$, but it will fail to capture some of the catastrophic failures of the very low shot count regime.

Average values of the shift $S$ and $E_\mathrm{proj}$ are obtained by a reblocking analysis\cite{Flyvbjerg1989} of data obtained in the steady-state regime of calculations, using pyblock.\cite{pyblock}

\begin{figure}
    \centering
    \includegraphics[width=0.45 \textwidth, trim = 0 0cm 0 1.5cm, clip]{H3+-binding-S2-pyblock.pdf}
    \includegraphics[width=0.45 \textwidth, trim = 0 0cm 0 2cm, clip]{H3+-binding-E2-pyblock.pdf}
    \caption{\small \justifying Shift and $E_\mathrm{proj}$ estimates of the correlation energy of linear H$_3^+$ over the range $r_\mathrm{HH} = 0.5 - 3.0 a_0$. Deviation from the true correlation energy increases as the number of shots is decreased, with deviations more pronounced in the projected energy and at larger bond lengths. This correlates with an increase in strong correlation and lowering of the overlap of the ground state wavefunction and the Hartree--Fock reference, $\braket{\phi_0|\Psi(\beta)}$.} 
    \label{fig:binding}
\end{figure}

\subsection{Ground State Calculations}

We begin our investigation into the effects of sampling noise on the projection algorithm by carrying out calculations using the circuit in \Cref{fig:circs} (a), sampled with $10^1 \mathrm{ to } 10^4$ shots for each Pauli term in the Hamiltonian. 

Representative results for H$_3^+$ at a bond length of $r = 2.0$ \r{A} are shown in \Cref{fig:shot_noise}. We note that, unsurprisingly, the noise in the energy estimators decreases as the number of shots is increased (see \Cref{tab:variance}). There is a qualitative change in behaviour in the projected energy using $N_\mathrm{shots} = 10$, with some samples deviating significantly from the expected mean. This is effectively due to the underestimation of the denominator in the expression of $E_\mathrm{proj}$ (see \Cref{eq:e_proj}), which leads to divergences in this quantity. There is therefore a lower bound on the number of shots needed to obtain a stable calculation, which is dependent on the overlap $\braket{\phi_0|\Psi(\beta)}$.


\begin{table}
    \centering
    \begin{tabular}{c|c|c|c|c|c}
      System &  $\braket{E}/E_\mathrm{h}$ & $\sigma(E)/E_\mathrm{h}$ & $\braket{S}/E_\mathrm{h}$ & $\sigma(S)/E_\mathrm{h}$&$E_\mathrm{FCI}/E_\mathrm{h}$\\
      \hline
         H$_3^+$& -0.068 & 0.001 & -0.0688&0.0005&-0.06803 \\
         H$_3$ & -0.110.7& 0.0009& -0.113& 0.001&-0.11011\\
         H$_4$ & -0.176& 0.002& -0.180& 0.003& -0.16701\\
         \hline
    \end{tabular}
    \caption{\small \justifying Mean and standard deviation of the shift and projected energy obtained from the deterministic propagation of stochastically rounded wavefunctions for systems of increasing size. }
    \label{tab:stoch_wfn}
\end{table}

\begin{figure}[h!]
    \centering
    \includegraphics[width=0.45 \textwidth, trim = 0 1cm 0 1.5cm, clip]{stoch_wfn_S.pdf}
    \includegraphics[width=0.45 \textwidth, trim = 0 1cm 0 2cm, clip]{stoch_wfn_E.pdf}
    \caption{\small \justifying Shift (top) and projected energy (bottom) as a function of imaginary time, obtained from a stochastically represented wavefunction for H$_3^+$ (blue), H$_3$ (orange) and H$_4$ (green) at $r = 1.5$ \r{A}. The noise in the shift is roughly system-independent, whereas the variance of $E_\mathrm{proj}$ increases with its absolute value.}  
    \label{fig:stoch_wfn}
\end{figure}

We compare performance at this bond length with conventional  VQE and PQE in \cref{tab:comp}. Exact VQE was carried out using the L-BFGS optimiser,\cite{Liu1989} while shot-based VQE used the COBYLA algorithm.\cite{Powell1994}  Better than VQE accuracy can be obtained with a 20-fold decrease in required shot count.

\begin{table}[]
    \centering
    \begin{tabular}{c|c|c}
         Method & Iterations & $\Delta E/\mathrm{m}E_\mathrm{h}$ \\
         \hline
         VQE &  44 & - \\
         VQE 10000 shots & 44 & 8\\
         PQE & 112 & - \\
         MC-PQE 100 shots & 250 & 0.4(1)\\
         \hline
    \end{tabular}
    \caption{\small \justifying Comparison of deterministic VQE and PQE results with shot-based VQE with $N_\mathrm{shots} = 10000$ and Monte Carlo estimated PQE results with $N_\mathrm{shots} = 100$.}
    \label{tab:comp}
\end{table}

Average values for the estimators $S$ and $E_\mathrm{proj}$ along the binding curve of H$_3^+$ are given in \Cref{fig:binding}. At large bond-lengths, we observe that divergences in the instantaneous values of the estimators due to undersampling translate into divergences of the overall estimator, with the projected energy more sensitive to this than the shift. This behaviour is not unexpected. As the bond length increases, the system becomes increasingly strongly correlated and therefore the weight of the reference determinant $\phi_0$ decreases, requiring more sampling to estimate it appropriately.

It is therefore the case that more strongly correlated systems will require more sampling than weakly correlated ones. This follows well-known trends in the difficulty of solving quantum chemical problems in a classical setting. We can however offer a series of ways to alleviate this issue. Firstly, in the absence of noise, the overlap $\braket{\phi_0|\Psi(\beta)}$ decreases monotonically with $\beta$. Therefore, it will be possible to use fewer shots at short imaginary times and increasing the shot count as the calculation progresses, thereby accelerating convergence to the steady state. For most Monte Carlo calculations however, the majority of the time is spent sampling this state, which cannot benefit from the early reduction in shots. 

\begin{table}
    \centering
    \begin{tabular}{c|c|c|c|c|c}
      $N_\mathrm{hamil}$ &  $\braket{E}/E_\mathrm{h}$ & $\sigma(E)/E_\mathrm{h}$ & $\braket{S}/E_\mathrm{h}$ & $\sigma(S)/E_\mathrm{h}$&$E_\mathrm{FCI}/E_\mathrm{h}$\\
      \hline
         2&-0.125&0.007&-115&3&\multirow{5}{*}{-0.11586}\\
         3& -0.121& 0.004& -0.116& 0.001&\\
         4 & -0.121& 0.004& -0.117& 0.001&\\
         5 & -0.115& 0.004& -0.1137& 0.0007&\\
         6 & -0.112 & 0.003 & -0.1133& 0.0007&\\
         \hline
    \end{tabular}
    \caption{\small \justifying Mean and standard deviation of the shift and projected energy obtained from linearised imaginary time-propagation of the H$_3^+$ wavefunction at $r=2.0$ \r{A}, using stochastically sampled hamiltonians with increasing numbers of terms selected per step.}
    \label{tab:hamil}
\end{table}


Secondly, different numbers of shots can be allocated to different quantities of interest. In this case, one could run a calculation with a small number of shots for the residuals, but a large number of shots for the reference overlap estimation. As the system size increases, this will represent a smaller fraction of the computational effort needed, thereby allowing calculations that are effectively as shot-efficient as ones with the lowest number of shots.

Finally, the standard computational chemistry approach would be to replace the Hartree--Fock reference wavefunction in the projected energy with a more complex, multi-determinental quantity $\ket{\Psi_\mathrm{T}}$ which is hoped to have larger overlap with the true ground state, giving a new projected energy of the form
\begin{equation}
    E'_\mathrm{proj} = \frac{\braket{\Psi_\mathrm{T}|\hat H| \Psi(\beta)}}{\braket{\Psi_\mathrm{T}|\Psi_(\beta)}} = \frac{\sum_{\mathbf{i} \in \Psi_\mathrm{T}}c_{\mathrm{T}\mathbf{i}}\braket{\phi_\mathbf{i}|\hat H|\Psi(\beta)}}{\sum_{\mathbf{i} \in \Psi_\mathrm{T}}c_{\mathrm{T}\mathbf{i}}\braket{\phi_\mathbf{i}|\Psi(\beta)}}
\end{equation}

On a classical machine, this would require a linear increase in computation, however in the current scheme, all terms in both the numerator and denominator are already computed as part of the propagation. 

\begin{figure}[h]
    \centering
    \includegraphics[width=0.45 \textwidth, trim = 0 1cm 0 1.5cm, clip]{H3+-S-hamil-2.0.pdf}\\
    \includegraphics[width=0.45 \textwidth, trim = 0 1cm 0 2cm, clip]{H3+-E-hamil-2.0.pdf}
    \caption{\small \justifying Shift (top) and projected energy (bottom) as a function of imaginary time obtained by imaginary-time propagation using a stochastically sampled Hamiltonian with only 2-6 Pauli groups selected at each time-step for H$_3^+$ at $r = 2.0$ \r{A}. The full H$_3^+$ Hamiltonian has 62 Pauli terms, split here into 25 commuting groups.}
    \label{fig:hamil}
\end{figure}

\begin{figure*}
    \centering
    \includegraphics[width=0.45 \textwidth, trim = 0 1cm 0 1.5cm, clip]{stoch_wfn_shots_S.pdf}
    \includegraphics[width=0.45 \textwidth, trim = 0 1cm 0 2cm, clip]{stoch_wfn_shotsE.pdf}
    \caption{\small \justifying Shift (4 left panels) and $E_\mathrm{proj}$ (4 right panels) in a simulation of H$_3^+$ at $r = 1.75$, using 100 or 1000 shots per residual estimation and shift damping parameters $\zeta = 1$ in the undamped case and $\zeta=0.1$ in the damped case. }
    \label{fig:shift_damping}
\end{figure*}

We move on to consider stochastic implementations of the projection algorithm. \Cref{fig:stoch_wfn} shows trajectories computed using deterministic projection on a stochastically rounded wavefunction for H$_3^+$, H$_3$ and H$_4$, with corresponding statistical properties given in \Cref{tab:stoch_wfn}. H$_2$ is not considered here as it only has one parameter in the wavefunction, so there is no rounding to be done. While the noise introduced into the shift by this procedure is roughly constant and system-independent, in itself a desirable property, the noise in $E_\mathrm{proj}$ scales linearly with the absolute magnitude of the correlation energy. This is unsurprising as if all amplitudes are rounded down, we expect the projected energy to be zero. Nevertheless, for all three systems we obtain stable estimators for the correlation energy. For the H$_3$ and H$_3^+$ systems, for which the UCCSD Ansatz is exact, the results are within error bars of the FCI value. For H$_4$, this approximation introduces an error of ~10 milliHartree in the correlation energy. 

\begin{figure}
    \centering
    \includegraphics[width=0.45 \textwidth, trim = 0 1cm 0 1.5cm, clip]{H3+-S-hamil-stoch-1000shots-1.75.pdf}
    \includegraphics[width=0.45 \textwidth, trim = 0 1cm 0 2cm, clip]{H3+-E-hamil-stoch-1000shots-1.75.pdf}
    \caption{\small \justifying Shift (top) and projected energy (bottom) as a function of imaginary time obtained by imaginary-time propagation using a stochastically sampled Hamiltonian with only 2-11 Pauli groups selected at each time-step, a stochastically rounded wavefunction and 1000 shots.} 
    \label{fig:full_stoch}
\end{figure}


\begin{figure}
    \centering
    \includegraphics[width=0.45 \textwidth, trim = 0 1cm 0 1.5cm, clip]{H3+-binding-S-full-stoch.pdf}
    \includegraphics[width=0.45 \textwidth, trim = 0 1cm 0 2cm, clip]{H3+-binding-E-full-stoch.pdf}
    \caption{\small \justifying Shift (top) and projected energy (bottom) estimators for the H$_3^+$ correlation energy, obtained using a sampled Hamiltonian with only 2-11 Pauli groups selected at each time-step, a stochastically rounded wavefunction and 1000 shots per residual.} 
    \label{fig:full_stoch_bind}
\end{figure}

\Cref{fig:hamil} shows the results of propagating the full wavefunction with a stochastically sampled Hamiltonian in which only $N_\mathrm{hamil}$ Pauli groups are included at each time-step. Corresponding values of the estimators are given in \Cref{tab:hamil}. Once again, as expected, variance decreases with the number of terms included and we note happily that propagation under the action of the stochastically sampled Hamiltonian does not introduce any aditional bias to the estimators. Noise in the shift fully depends on the details of the stochastic trajectory. For estimation of $E_\mathrm{proj}$ however we use the full Hamiltonian at each step, which helps reduce noise significantly. As this is essentially equivalent to one residual computation, the relative cost of employing the full Hamiltonian at this stage decays with system size. However, the projected energy could also be estimated using the sampled Hamiltonian at each step, decreasing computational overhead at the cost of increased sampling noise.


Considering the data presented above, we are satisfied that shot-based simulation, stochastic wavefunction rounding and Hamiltonian sampling are all compatible with imaginary-time evolution and do not introduce significant biases into the results of a simulation. They do however all increase the noise present, requiring, much like classical QMC methods, a relatively long period of steady-state propagation to converge the desired quantities to the intended accuracy.


There are ways to reduce the noise in the simulation which do not require increased sampling. One such method is \textit{shift damping},\cite{Booth2009} which entails decreasing the parameter $\zeta$ in \Cref{eq:shift}. This makes the shift less responsive to changes in population, thereby directly decreasing its variance. As the value of the shift affects the continued propagation of the wavefunction, we expect that, once the steady-state has been reached, less variability of the shift will translate to less variability of the wavefunction and therefore one might expect less noise in the projected energy as well. \Cref{fig:shift_damping} compares the shift and projected energy in calculations with different numbers of shots and different shift damping, as well as a stochastic representation of the wavefunction. While decreasing $\zeta$ is very effective in reducing shift noise ($\sigma_\mathrm{damped}/\sigma =$ 18\% (1000 shots), 20\% (100 shots)), we see no effect in the projected energy ($\sigma_\mathrm{damped}/\sigma =$ 101\% (1000 shots), 102\% (100 shots)).

Finally, in \Cref{fig:full_stoch} we combine shot-based simulation with stochastic wavefunction rounding and hamiltonian sampling. It is clear that simulations remain well behaved even with multiple sources of noise. We also note that the systematic increase in error due to stochastic rounding of the wavefunction appears to be partially cancelled out by other noise. Binding curves obtained with the fully stochastic algorithm are given in \Cref{fig:full_stoch_bind}. As we noted in the case of finite sampling noise, more highly correlated regimes require better Hamiltonian sampling to achieve the same level of accuracy.

\begin{table*}
    \centering
    \begin{tabular}{l||c|c|c|c||c|c|c|c||c|c|c|c}
      \multirow{2}{*}{System} &  \multicolumn{2}{c|}{MAE ($\sigma$)} & \multicolumn{2}{c||}{NPE ($\sigma$)} & \multicolumn{2}{c|}{MAE ($\sigma$, $N_h = 10$)} & \multicolumn{2}{c||}{NPE ($\sigma$, $N_h = 10$)} & \multicolumn{2}{c|}{MAE ($\sigma$, $N_h = 20$)} &\multicolumn{2}{c}{NPE ($\sigma$, $N_h = 20$)}\\
      \cline{2-13}
      & S & E& S & E& S & E& S & E& S & E& S & E\\
      \hline
        LiH & 0.5 & 0.3 & 0.6 & 0.6 & \hspace*{5mm}2\hspace*{5mm} & 1& \hspace*{5mm}4\hspace*{5mm} &1.6&\hspace*{5mm}2\hspace*{5mm} &0.7 &\hspace*{4mm}2.7\hspace*{4mm}&1\\
         BeH$_2$ & 0.6 & 0.9 & 0.5 & 0.9 & 2 & 3 & 5 & 6 &2&3&4&4 \\
         H$_2$O & 6 & 4 & 6& 6& 42 &76 & 53& 72&9&5& 10& 5 \\
         HF& 9 & 4& 14 & 4 & 10& 13& 18 & 20 & 6 & 10 & 6 &14\\
         \hline
    \end{tabular}
    \caption{\small \justifying Maximum absolute error (MAE) and non-parallelity error (NPE) in milliHartree for Monte Carlo PQE in some frozen-core second row hydrides, compared to fully deterministic PQE. These systems have 125 (LiH and HF) and 313 (H$_2$O and BeH$_2$) Hamiltonian groups respectively. Results in all columns have Gaussian noise with $\sigma = 0.001$ added to simulate sampling noise and results are obtained from runs of 2000 imaginary time-steps $\delta \beta = 0.2$. The first set of results (columns 1-4) have no additional stochasticity. The second (5-8) and third (9-12) sets use stochastic wavefunctions and 10 or 20 sampled Hamiltonian terms per step respectively. The error converges towards the pure sampling error with few sampled Hamiltonian groups. For all methods, errors could be further decreased by using longer propagation times.}
    \label{tab:hydrides}
\end{table*}

We apply this stochastic PQE algorithm to first row mono- and dihydrides over a range of geometries, with results summarised in \Cref{tab:hydrides}. The cases studied include the dissociation of LiH and HF, the symmetric stretch of H$_2$O at its equilibrium bond-angle and the $C_{2v}$ insertion of Be into the H$_2$ bond to form BeH$_2$.\cite{Purvis1983}
Geometries for HF and H$_2$O match those from \cite{Das2010} and the STO-3G basis set was used for all molecules. When using the same calculation parameters for all systems, it is unsurprising that errors are larger for more complex molecules. However, we note that by stochastically rounding the wavefunction and only sampling 20 Hamiltonian groups for each system, similar errors to the pure finite sampling case are observed. In all cases, the results correspond to relatively short propagation times and errors could be reduced by running the calculations further.

\subsection{Excited States}

As described in Sec II D, the algorithms above can also be applied to excited states by employing the modified Hamiltonian in \Cref{eq:fs_ham}. The results of example folded-spectrum calculations  performed using exact statevector simulations for H$_3^+$ are given in \Cref{fig:fs_h3+}. In order to easily converge to the desired state, the calculations are initialised at the first geometry (shortest $r_\mathrm{HH}$) with different reference determinants and values of $\omega$ close to their energies. 
States are then tracked along the binding curve by using the parameters and energies of the previous point to initialise the wavefunction and $\omega$ value for the current point. This allows smooth tracking of the states over the whole range of geometries considered.

\begin{figure}
    \centering
    \includegraphics[width=0.5\textwidth]{H3+-fs.pdf}
    \caption{\small \justifying Statevector-simulated folded-spectrum energies for the H$_3^+$ molecule in the STO-3G basis set, using a projective eigenvalue solver. The Slater determinants used as initial references for each excited states are, in increasing energy order of the resulting state: $\ket{01}$, $\ket{03}$, $\ket{12}$, $\ket{23}$, $\ket{05}$, $\ket{14}$, $\ket{25}$, $\ket{34}$, $\ket{45}$.}
    \label{fig:fs_h3+}
\end{figure}

We investigate here introducing noise in these simulations, first in the form of pure sampling noise due to shot-based simulation. Results obtained using Gaussian noise with $\sigma = 0.01$ for  H$_3^+$ are shown in \Cref{fig:fs_reblock}. We note that, even with relatively realistic finite sampling noise levels, convergence is still very good, with obtained states closely following their FCI  counterparts.

\begin{figure}
    \centering
    \includegraphics[width=0.5\textwidth]{H3+-100shots.pdf}
    \caption{\small \justifying Folded-spectrum energies for the H$_3^+$ molecule in the STO-3G basis set, using a projective eigensolver and Gaussian noise with $\sigma = 0.01$.}
    \label{fig:fs_reblock}
\end{figure}

However, not all instances of noise are as benign in this scenario, as can be seen in \Cref{fig:fs_stoch_wfn}, which shows results for the totally symmetric states of H$_3^+$ obtained with stochastically rounded wavefunctions. In this case, particularly for the highly excited states, the rounding of the wavefunction leads to unphysical behaviour. The third excited state is most strongly affected, at points tending towards converging onto the highest excited state instead. One must note that because previous values of the energies are used as $\omega$ values at subsequent points, once a sufficiently large deviation occurs it is likely to remain or be amplified. This could be avoided by running independent calculations at each geometry, but a new problem of estimating \textit{a priori} a good initial wavefunction and $\omega$ is introduced.
\begin{figure}
    \centering
    \includegraphics[width=0.5\textwidth]{fs-shots-reblock.pdf}
    \caption{\small \justifying Folded-spectrum energies for the totally symmetric states of the H$_3^+$ molecule in the STO-3G basis set, using a projective eigenvalue solver and stochastically rounded wavefunctions at each step.}
    \label{fig:fs_stoch_wfn}
\end{figure}

The main limitation to using the folded spectrum method is the increase in terms in the squared Hamiltonian compared to the original operator. For non-trivial systems, this leads to a potentially insurmountable increase in required computation, even with efficient Hamiltonian term grouping methods. However, the Monte Carlo methods described above provide an effective means of sampling the folded spectrum operator. We find (see \Cref{fig:fs_hamil}) that the same number of samples can be used as in a standard ground state calculations to generate good quality excited state estimates. This significantly expands on the range of systems for which folded spectrum methods are tractable.

\begin{figure}
    \centering
    \includegraphics[width=0.5\textwidth]{H3+-100shots-hamil.pdf}
    \caption{\small \justifying Folded-spectrum energies for the H$_3^+$ molecule in the STO-3G basis set, using a projective eigensolver with $\sigma = 0.01$ Gaussian noise and 10 Hamiltonian terms sampled at each step.}
    \label{fig:fs_hamil}
\end{figure}

\section{Conclusion}
In this paper, we present a Quantum Monte Carlo-informed approach to projective quantum eigensolver algorithms. We begin by noting that rather than attempting to obtain the true ground state of a system to within the desired accuracy directly, once propagation has converged, we can obtain highly accurate estimates by classically averaging over samples obtained at different times over the course of imaginary time propagation. This allows us to use low shot counts without compromising final accuracy, requiring up to 50 times fewer samples than equivalent conventional VQE calculations.

We go further to decrease the quantum computational overhead, through stochastic rounding of the wavefunction and probabilistic sampling of the Hamiltonian. These procedures lead to shallower circuits and fewer shots required per step, but also generate noisier projectors.  However, we observe that the Monte Carlo energy estimators employed are generally resilient to this additional noise, generating comparable accuracy to the deterministic projector, provided the resulting truncation at each step is not too drastic. In practice, for second row hydrides we observe that as few as 10\% of the Hamiltonian terms are required.

Finally, we expand this method to excited state calculations, by employing the folded-spectrum approach, in which the Hamiltonian is replaced by a variance operator.  The squared operator has significantly more terms than the underlying Hamiltonian, making this method traditionally more expensive than conventional eigensolvers. However, by sampling the operator we find that we can reduce the cost to that of an equivalent ground state calculation.

The fundamental reason these stochastic approximations and sampling methods work is that, given a wavefunction that is a steady-state of some operator, they induce fluctuations about this wavefunction that may be effectively averaged over. Therefore, their applicability is not limited to the projective quantum eigensolver and we believe they may be helpful in any quantum algorithm which aims to converge to a some well-defined final wavefunction.

\section*{Acknowledgements}
The author thanks Dr. Alex Thom for useful discussions and gratefully acknowledges Peterhouse, University of Cambridge for supporting this work through a Research Fellowship.

\bibliography{qc_qmc}
\end{document}
